\subsection{RQ1: How Can FPV Quadcopter Flight Physics Be Simulated?}

% TODO: answer

% FPV quadcopter flight physics can be effectively simulated using a simplified control model based on proportional-derivative (PD) control, particularly when the goal is to support rapid prototyping and beginner-friendly interaction in XR environments. 

% While existing simulators commonly employ cascaded PID controllers to achieve high physical fidelity, the findings of this study suggest that full PID complexity is not strictly necessary for early-stage learning and design exploration. Instead, a reduced model combined with self-leveling constraints provides a more accessible and stable interaction experience, especially when integrated with human-AI collaborative assistance.

This study adopts Yue's Ultimate Drone physics model \parencite{yuetility2022fpvphysics} as a reference implementation for simulating first-person view (FPV) quadcopter flight dynamics. The model employs a proportional-derivative (PD) control approach for flight stabilization, which is widely used in modern quadcopter systems.

In addition, a mapping strategy was developed to translate the controller profile of the Meta Quest into a simulated drone transmitter, thereby enabling intuitive user interaction within the virtual environment.

To ensure safe operation during simulation, the environment incorporates collider-based boundaries that constrain the quadcopter within a predefined area.

\subsubsection{Comparison With Existing Simulators}

FPV simulators (e.g., Liftoff, Uncrashed) primarily adopt cascaded proportional-integral-derivative (PID) control models to replicate real-world flight controllers. These systems typically support multiple controller configurations and provide interfaces for PID parameter tuning. However, such configurations can be difficult for beginners to set up with their own controllers.

Although most simulators provide a variety of selectable environments, they are primarily designed for flat-panel displays and do not offer the immersive learning enabled by extended reality (XR) environments.

Commercial FPV simulators commonly employ Acro (acrobat) mode to achieve realistic flight behavior. In contrast, this study focuses on a self-leveling flight mode, which constrains the drone's body angle and reduces the likelihood of crashes, thereby minimizing user frustration during early-stage learning.

Unlike existing simulators, the proposed system incorporates a human-AI collaboration mechanism that provides calculated control suggestions to support user decision-making.

Although Acro mode remains important for future investigation due to its higher realism, it also introduces a greater likelihood of crashes. This raises several design considerations, including mechanisms for crash recovery, feedback systems for crash detection, and strategies for enabling AI agents to avoid collisions. From a game design perspective, it is also necessary to evaluate whether and how penalties should be applied following crashes.
